% =============================================================================
% Chapter 6 -- Project Ethics
% =============================================================================
\chapter{Project Ethics}
\label{ch:ethics}

The COMP390 handbook requires every dissertation to include an
explicit ethics chapter, regardless of whether human participants
were involved. This chapter records the ethics-relevant decisions
that were made during the project.

\section{Data and Participants}

This dissertation studies the behaviour of reinforcement learning
agents trained in a simulated game environment. There are no human
participants, no personal data is processed, and no third-party
data is collected or redistributed beyond what is bundled with the
public COOM and ViZDoom releases. The data category and
participant category recorded on the Statement of Ethical
Compliance reflect this. No ethics review board approval was
required, and none was sought.

\section{Wider Implications}
\label{sec:wider}

Continual reinforcement learning is a dual-use technology.
Improvements in the field could enable cheaper on-device
adaptation of beneficial autonomous systems --- domestic
assistive robots, accessibility tools, energy-management
controllers --- but the same techniques apply equally to
autonomous weapons platforms, surveillance agents and other
applications that the BCS code of conduct identifies as ethically
fraught.

The contribution of this dissertation is, on balance, a
\emph{deflationary} one: the negative-result framing argues
\emph{against} premature deployment of regularisation-based CL
under the compute conditions a small operator is likely to face.
That direction is consistent with the precautionary stance taken
by both the BCS code of conduct, which requires members to act in
the public interest and to give due regard to the wider impact
of their work, and the ACM code of ethics, which similarly
emphasises avoiding harm and acknowledging the limits of one's
findings. By reporting an inversion of the published ranking at
reduced compute, the dissertation hopes to lower the probability
that a deployer mistakes a benchmark headline for evidence of
real-world readiness.

\section{Use of Generative AI Tools}
\label{sec:genai}

The COMP390 handbook requires explicit disclosure of any
generative AI tool use. During the project, generative AI tools
were used in three bounded ways:
\begin{itemize}
    \item \textbf{Editorial assistance} on the prose of this
        dissertation: rephrasing draft passages and shortening
        sentences. All claims, all numbers and the argumentative
        structure of every chapter are the author's own.
    \item \textbf{Boilerplate code generation} for plotting and
        log-loading utilities under \texttt{results/plotting/}
        and \texttt{paper/scripts/}. Generated code was reviewed,
        tested against the recorded \texttt{progress.tsv} files,
        and modified before use.
    \item \textbf{Literature search aids}, used to surface
        candidate references that were then read and verified
        against their primary sources before being cited.
\end{itemize}
No generative tool was used to design the experiments, to
choose hyperparameters, to interpret the results, or to write
the self-reflection in \Cref{ch:bcs}. No generated text was
included in the final dissertation without authorial review and
revision.
